\chapter*{课程说明}
\addcontentsline{toc}{chapter}{课程说明}
这门课程讨论三个相互连接的问题：一个优化问题在什么条件下有可验证的最优解；一个迭代方法为什么下降、收敛或失稳；数据、参数化和更新规则怎样共同决定神经网络的训练动力学。前两讲以凸优化、对偶与KKT条件建立分析语言，后三讲沿着随机梯度、预条件、矩阵更新与宽度缩放展开。

讲义服务于五次课，共16学时，其中4学时为实验。先修知识为微积分、线性代数与基础概率；实验需要能够阅读Python数组运算。每讲既有课堂核心，也有供研究生课后复习的补充推导。无需在有限课时内逐页讲完所有选读内容。

\begin{center}
\begin{tabularx}{\textwidth}{cXrr}
\toprule
讲次 & 主题 & 理论 & 实验\\
\midrule
1 & 凸优化与下降方法 & 4 & 0\\
2 & 对偶、KKT与约束优化 & 3 & 0\\
3 & 随机优化、动量与非凸收敛 & 2 & 1\\
4 & 预条件、自适应优化与矩阵更新 & 1.5 & 1.5\\
5 & 参数化、NTK与跨宽度迁移 & 1.5 & 1.5\\
\bottomrule
\end{tabularx}
\end{center}

\section*{如何使用本讲义}
阅读时先明确优化变量与目标，再核对定理的条件。每个算法应回答四个问题：输入与状态是什么，更新式是什么，停止指标是什么，保证依赖什么假设。遇到缩放公式时，还需写出实际参与前向计算的参数和代码中被更新的变量。

\begin{coursebox}{四种结论的身份}
\begin{description}[style=nextline,leftmargin=1em]
\item[精确结论] 在列出的模型与条件下可以证明，包括代数恒等式、KKT充分性和二次模型稳定条件。
\item[近似模型] 例如梯度流、局部线性化、小步长扩散近似。结论的使用必须保留近似范围。
\item[经验观察] 例如有限宽度超参数迁移、特定训练设置的更新RMS和稳定性边界。需要可复现实验支撑。
\item[研究问题] 当前材料不能推出的一般结论，例如仅由一个曲率标量判断所有模型的泛化。
\end{description}
\end{coursebox}

\section*{统一记号}
向量默认是列向量。$\norm{x}_2$是欧氏范数，$\norm{A}_2$是矩阵谱范数，$\norm{A}_F$是Frobenius范数。$n$通常为样本数，$d$为输入维数，$m$为隐藏宽度，$B$为batch大小，$T$为训练步数。$L$在收敛分析中表示梯度Lipschitz常数；若某个例子使用其他局部记号，会在原处重新说明。$F$通常为总体或经验目标，$\ell$为单样本损失。

学习问题采用平均损失时显式保留$1/n$或$1/B$。二次模型中若使用未归一化平方损失，其步长尺度会相应改变。连续时间用上点表示导数，离散迭代使用下标$t$；连续梯度流恒等式不能自动视为任意有限步长的结论。

\section*{实验与课后题}
实验代码位于同目录的\texttt{experiments/}，交互式入口位于\texttt{notebooks/}。所有数据均由公开的合成规则本地生成，无需账户或外部数据下载。PyTorch实验默认自动选择可用设备，也支持显式指定CPU或CUDA。比较优化器时固定数据划分、基础模型和调参预算；比较宽度时固定训练步数、batch与评价探针，报告初始化尺度。

每讲末尾提供习题与解答要点。先独立推导，再核对答案。课程考核建议同时检查假设、计算与解释：只写出最终公式不足以说明其适用性。
